\chapter{MACVIO: Learning-Based Visual-Inertial Odometry with Cross-Modal Uncertainty Fusion} \label{chapter:macvio}
\fancyhead[LO]{\makebox[\textwidth][r]{Chapter 7. MACVIO}}

\section{Introduction}

This chapter presents MACVIO (Metrics-Aware Cross-Modal Visual-Inertial Odometry), which extends the single-sensor uncertainty modeling from MAC-VO and AirIMU into a unified visual-inertial estimator.
MACVIO builds directly on three foundations established earlier in this dissertation: metrics-aware visual covariance from MAC-VO, learned inertial uncertainty from AirIMU and AirIO, and robust initialization and online calibration from MAC-I$^2$.
The central idea is to fuse uncertainty information from both modalities and adapt fusion behavior online so that VIO remains accurate and stable under lighting degradation, aggressive motion, dynamic scenes, and sensor-quality variation.

Conventional visual-inertial systems rely on fixed fusion weights and manually tuned noise parameters.
These settings are brittle under distribution shift because visual and inertial reliability vary with environment and motion context.
MACVIO addresses this limitation by learning cross-modal uncertainty correlation and converting it into adaptive fusion weights during optimization.

The motivation for MACVIO is practical deployment.
In real robotic systems, sensing quality changes rapidly: cameras degrade in low light or high-speed motion, while inertial signals degrade under vibration, bias drift, and unmodeled dynamics.
A fixed-weight fusion strategy cannot respond to these shifts and often fails exactly when robustness is most needed.
Therefore, this chapter introduces uncertainty-aware cross-modal control as the layer between modality-specific uncertainty prediction and full SLAM integration.

From a thesis-structure perspective, MACVIO is the bridge chapter in Part~III.
Parts~I and II establish reliable uncertainty for individual modalities.
MACVIO transforms these single-modal signals into a joint decision mechanism that determines how much each modality should influence state estimation at each time step.
This mechanism improves transfer across robots and environments while reducing manual retuning effort.

\section{Method}
\label{sec:macvio:method}

\subsection{Basic Idea}

The basic idea of MACVIO is to treat visual and inertial uncertainty as complementary signals and fuse them with a learned reliability model.
Given visual measurements $\mathcal{Z}_{1:T}^{v}$ and inertial measurements $\mathcal{Z}_{1:T}^{i}$, we estimate state trajectory $\mathcal{X}_{1:T}$ by minimizing:
\begin{equation}
\label{eq:macvio:objective}
\mathcal{X}_{1:T}^{*}
=
\argmin_{\mathcal{X}_{1:T}}
\sum_{k} w_k^{v}\, \left\| r_k^{v}(\mathcal{X}_{1:T}) \right\|_{{\Sigma}_k^{v}}^2
+
\sum_{k} w_k^{i}\, \left\| r_k^{i}(\mathcal{X}_{1:T}) \right\|_{{\Sigma}_k^{i}}^2 ,
\end{equation}
where ${\Sigma}_k^{v}$ is provided by the visual uncertainty model (MAC-VO), ${\Sigma}_k^{i}$ is provided by the inertial uncertainty model (AirIMU/AirIO), and $(w_k^{v}, w_k^{i})$ are adaptive fusion weights.

\subsection{Cross-Modal Reliability Learning}

MACVIO learns a cross-modal reliability state from uncertainty and context:
\begin{equation}
\label{eq:macvio:state}
\mathbf{s}_k = \phi\!\left({\Sigma}_k^{v}, {\Sigma}_k^{i}, \mathbf{c}_k\right),
\end{equation}
where $\mathbf{c}_k$ includes motion and feature-quality context (for example, image degradation and inertial excitation level).
The reliability state is mapped to adaptive weights by:
\begin{equation}
\label{eq:macvio:weights}
\left(w_k^{v}, w_k^{i}\right) = g(\mathbf{s}_k), \quad w_k^{v}, w_k^{i}\in[0,1].
\end{equation}
This mechanism increases trust in the currently reliable modality and suppresses degraded measurements.

\subsection{How Previous Chapters Are Fused}

MACVIO fuses outputs from prior chapters in a single VIO loop:
\begin{itemize}
    \item \textbf{From MAC-VO:} metrics-aware visual covariance for correspondence filtering and residual weighting.
    \item \textbf{From AirIMU and AirIO:} learned inertial covariance and observability-aware inertial representation for high-dynamic motion.
    \item \textbf{From MAC-I$^2$:} robust visual-inertial initialization and online calibration priors.
\end{itemize}
By combining these modules, MACVIO provides better accuracy than fixed-weight fusion, especially in scenes where one modality temporarily degrades.

\subsection{Training and Deployment}

The training strategy has two stages.
First, pretrained visual and inertial uncertainty modules are loaded from MAC-VO and AirIMU to provide stable covariance priors.
Second, the cross-modal fusion controller is trained to predict adaptive weights that minimize trajectory error and consistency loss.
During deployment, MACVIO updates weights online without per-platform retuning, aligning with the thesis goal of minimal additional human engineering.
